\section{The printed method}\label{sec:printed}

Throughout this section $\gamma\ge\gamma_0$.

\subsection*{Test field}
Let $\varphi_0(\theta):=\int_0^\theta(p-\pbar)(\theta')\,d\theta'$, which is periodic, and $\varphi:=\chi_1(r)\varphi_0(\theta)$, where $\chi_1\equiv1$ near $r=1$ (so $\chi_1'(1)=0$) and $\chi_1\equiv0$ near $\partial R$.
Then $w:=\curl\varphi$ has components $w_r=\partial_\theta\varphi/r$ and $w_\theta=-\partial_r\varphi$, and on $\Gamma$
\[
w=(p-\pbar)e_r=-(p-\pbar)n .
\]
Let $\Pi_{k+1}$ be the $C^1$ piecewise-$P_{k+1}$ interpolant of Morgan and Scott \cite{MorganScott1975} (for $k=4$, the Argyris interpolant \cite{Argyris1968}), which exists because $k+1\ge5$; its constants depend on the angles, hence on (M0) and (M2). Since $\varphi\equiv0$ near $\partial R$ and $h\le h_\ast$, $v_\varphi:=\curl\Pi_{k+1}\varphi\in\Zh$, and $\norm{v_\varphi-w}_{W^{1,\infty}}\le Ch^k$.

\begin{theorem}[Energy norm: exact rate $1/2$]\label{thm:A}
The error satisfies the upper bound
\[
\tnorm{\ut-u_h}\le C\Bigl[(\hG/\gamma)^{1/2}G+(1+\gamma^{1/2})\hG^{3/2}+\eps\Bigr].
\]
If $G>0$ and $h\le h_0(G)$, with $h_0$ independent of $\gamma\ge\gamma_0$, it also satisfies the lower bound
\[
\tnorm{\ut-u_h}\ge(2M)^{-1}(\hG/\gamma)^{1/2}G-C\Bigl[(1+\gamma^{1/2})\hG^{3/2}+\eps\Bigr].
\]
In Scott's variables, $(\hG/\gamma)^{1/2}=(h/\mu)^{1/2}$.
\end{theorem}

\begin{proof}[Proof of the upper bound]
For $z\in\Wh$, $e_h:=z-u_h\in\Zh$. By Corollary~\ref{cor:GSerr} and Lemma~\ref{lem:coercive},
\[
c_1\tnorm{e_h}^2\le N_h(z-\ut,e_h)+\Rc(e_h)+\Gc(e_h).
\]
Using the Jacobian bound of Lemma~\ref{lem:chords}, $|\Rc(e_h)|\le(G+C\hG^2)(h/\mu)^{1/2}\tnorm{e_h}$. Then apply Lemmas~\ref{lem:Gbound} and~\ref{lem:approx}.
\end{proof}

\begin{proof}[Proof of the lower bound]
\begin{enumerate}
\item We have $v_\varphi\cdot n_e=-(p-\pbar)(x^\ast)\,n(x^\ast)\cdot n_e+O(h^k)$, $\Rc=-\ip{\pt-\pbar}{v\cdot n}$ and $n(x^\ast)\cdot n_e=1+O(s^2)$. Hence
\[
\Rc(v_\varphi)=\ip{(\pt-\pbar)\,(p-\pbar)(x^\ast)}{n(x^\ast)\cdot n_e}+O(h^k)=G^2+O(\hG^2+h^k).
\]
\item $\tnorm{v_\varphi}^2\le(\mu/h)(G+C\hG^2+Ch^k)^2+C$.
\item By continuity, $M\tnorm{\ut-u_h}\ge(|\Rc(v_\varphi)|-|\Gc(v_\varphi)|)/\tnorm{v_\varphi}$, and Lemma~\ref{lem:Gbound} bounds $\Gc$.
\item $|\Rc|/\tnorm{v_\varphi}\ge(h/\mu)^{1/2}G\bigl(1-C(\hG^2+h^k)/G^2-C(\hG^2+h^k)/G-C(\hG/\gamma_0)^{1/2}/G\bigr)$, which is at least $\frac12(h/\mu)^{1/2}G$ for $h\le h_0(G)$.\qedhere
\end{enumerate}
\end{proof}

The lower bound is carried by the boundary parts of $\tnorm\cdot$. By Theorem~\ref{thm:C}, $\norm{\nabla(\ut-u_h)}$ is only $O(h/\mu)$.

\begin{corollary}\label{cor:Aprime}
If $G>0$, $\gamma\hG\ll G$ and $\eps\ll(\hG/\gamma)^{1/2}G$, then $\tnorm{\ut-u_h}\asymp(h/\mu)^{1/2}$. The mesh-dependent energy rate is exactly $1/2$.
\end{corollary}

\begin{lemma}[Normal-load cancellation]\label{lem:cancel}
Let $S$ be edgewise $W^{1,\infty}$ on $\Gh$, continuous at the vertices of $\Gh$, with $|S\cdot\tau_e|\le C\hG$ on every $e\in\EG$. (For instance, $S$ smooth with $S\cdot\tau=0$ on $\Gamma$.) Then for $v\in\Zh$,
\[
|\ip{S}{\dn v}|\le C\bigl(\norm{v}_{\Gh}+\hG^{1/2}\norm{\nabla v}\bigr)\le C\norm{\nabla v}.
\]
\end{lemma}

\begin{proof}
\begin{enumerate}
\item \emph{Tangential part.} By Lemma~\ref{lem:inverse} and (M0), $|\ip{(S\cdot\tau)\tau}{\dn v}|\le C\hG\norm{\dn v}_{\Gh}\le C\hG^{1/2}\norm{\nabla v}$.
\item \emph{Normal part.} On $e$, $v$ is a polynomial on $T_e$ with $\div v=0$, so $\dn v\cdot n_e=-\dt(v\cdot\tau_e)$. With $\sigma_e:=S\cdot n_e$,
\[
\int_e\sigma_e\,\dn v\cdot n_e=\int_e\dt\sigma_e\,(v\cdot\tau_e)-\bigl[\sigma_e\,v\cdot\tau_e\bigr]_{\partial e}.
\]
\item At a polygon vertex $x$ shared by $e$ and $e'$, the endpoint terms combine to $v(x)\cdot(\sigma_{e'}\tau_{e'}-\sigma_e\tau_e)$, of size at most $C(|e|+|e'|)|v(x)|$.
\item $\sum_e|e|\norm{v}_{L^\infty(e)}\le C\sum_e|e|^{1/2}\norm{v}_{L^2(e)}\le C|\Gh|^{1/2}\norm{v}_{\Gh}$.
\item Finish with Lemma~\ref{lem:korn}.\qedhere
\end{enumerate}
\end{proof}

Numerically (Section~\ref{sec:numerics}), the dual norm of $v\mapsto\ip{S}{\dn v}$ stays bounded on $\Zh$ for normal $S$. It grows like $\hG^{-1/2}$ on $\VR$ (no divergence constraint), and also on $\Zh$ for tangential $S$. Incompressibility is exactly what the proof uses.

\begin{theorem}[$H^1$ seminorm: rate $1$]\label{thm:C}
With $C_p:=C(1+\norm{p-\pbar}_{W^{1,\infty}(\Gamma)})$,
\[
\norm{\nabla(\ut-u_h)}\le C_p\,(\hG/\gamma)+C\Bigl[(1+\gamma^{1/2})\hG^{3/2}+\eps\Bigr].
\]
\end{theorem}

\begin{proof}
\begin{enumerate}
\item \emph{Split.} For $z\in\Wh$, write $z-u_h=\delta_R+\delta_G$, where $\delta_R,\delta_G\in\Zh$ solve
\[
N_h(\delta_R,v)=\Rc(v),\qquad N_h(\delta_G,v)=\Gc(v)-N_h(\ut-z,v)\qquad\forall v\in\Zh.
\]
Lax--Milgram on $\Zh$ (Lemma~\ref{lem:coercive}) gives existence and uniqueness.
\item \emph{Geometric part.} $\norm{\nabla\delta_G}\le\tnorm{\delta_G}\le c_1^{-1}\bigl(C(1+\gamma^{1/2})\hG^{3/2}+ME_h\bigr)$.
\item \emph{Ansatz.} Heuristically $\delta_R\approx-(h/\mu)(p-\pbar)n=(h/\mu)w$ on $\Gh$. Set $S_\varphi:=v_\varphi\in\Zh$ and $r_h:=\delta_R-(h/\mu)S_\varphi\in\Zh$. Then $N_h(r_h,v)=m(v)+\ell(v)$ with
\begin{align*}
m(v)&:=\Rc(v)-\ip{S_\varphi}{v}=-\ip{(\pt-\pbar)n+S_\varphi}{v},\\
\ell(v)&:=-(h/\mu)\bigl[a_h(S_\varphi,v)-\ip{\dn S_\varphi}{v}-\ip{S_\varphi}{\dn v}\bigr].
\end{align*}
\item \emph{Bound on $m$.} On $e$,
\[
(\pt-\pbar)n_e+S_\varphi=(p-\pbar)(x^\ast)(n_e-n(x^\ast))+O(\hG^2)+O(h^k)=(p-\pbar)(m_e^\ast)\,s\,\tau_e+O(\hG^2)+O(h^k).
\]
The leading term is odd in $s$ with a constant coefficient. The argument of Lemma~\ref{lem:transposed} handles it, and the $O(h^k)$ remainder contributes $Ch^k\norm{v}_{L^1(\Gh)}$. Hence $|m(v)|\le C_p(\hG^{3/2}+h^k)\norm{\nabla v}$.
\item \emph{Bound on $\ell$.} $|a_h(S_\varphi,v)|+|\ip{\dn S_\varphi}{v}|\le C_p\norm{\nabla v}$. By Lemma~\ref{lem:cancel} applied to $S=w$,
\[
|\ip{S_\varphi}{\dn v}|\le|\ip{w}{\dn v}|+|\ip{S_\varphi-w}{\dn v}|\le C_p\norm{\nabla v}+Ch^k\hG^{-1/2}\norm{\nabla v}.
\]
So $|\ell(v)|\le(h/\mu)(C_p+Ch^k\hG^{-1/2})\norm{\nabla v}$.
\item \emph{Close.} Testing with $v=r_h$,
\[
c_1\tnorm{r_h}^2\le\bigl[C_p(h/\mu+\hG^{3/2}+h^k)+C(h/\mu)h^k\hG^{-1/2}\bigr]\norm{\nabla r_h}.
\]
The last two terms are at most $C\eps$, since $(h/\mu)h^k\hG^{-1/2}=\gamma^{-1}h^k\hG^{1/2}\le Ch^k$. Hence
\[
\norm{\nabla\delta_R}\le(h/\mu)\norm{\nabla S_\varphi}+\norm{\nabla r_h}\le C_p(\hG/\gamma)+C(\hG^{3/2}+\eps).\qedhere
\]
\end{enumerate}
\end{proof}

\begin{remark}[Why rate $1$ and not $1/2$]
Without incompressibility, Lemma~\ref{lem:cancel} would give only $C\hG^{-1/2}$. The $H^1$ error would then be $(h/\mu)\hG^{-1/2}$: rate $1/2$, with element-scale oscillation in the trace. The inconsistency $-pn$ is purely normal, and $\div v=0$ converts $\dn v\cdot n$ into a tangential derivative. That is why $\GS$ still converges at rate $1$ in $H^1$.
\end{remark}

\begin{theorem}[Boundary slip; $H^1$ lower bound]\label{thm:B}
Let $\gamma\ge\gamma_0$ and $G>0$, and put $X:=C_p(\hG/\gamma)+(1+\gamma^{1/2})\hG^{3/2}+\eps$, the right-hand side of Theorem~\ref{thm:C}. Suppose that, for a fixed small $\delta$,
\[
X\le\delta\,G\min(G,\hG^{1/2}),\qquad \hG/\gamma\le\delta G,\qquad \hG\le h_0(G).
\]
All three hold for fixed $\gamma$ and bounded $\rho$ as $\hG\to0$. Then
\begin{enumerate}[label=(\roman*)]
\item $c(\hG/\gamma)G\le\norm{\ut-u_h}_{L^2(\Gh)}\le C(\hG/\gamma)^{1/2}\tnorm{\ut-u_h}$;
\item $\norm{\nabla(\ut-u_h)}\ge c(\hG/\gamma)G-Ch^{k+1/2}$.
\end{enumerate}
Combined with Theorem~\ref{thm:C}, $\norm{\nabla(\ut-u_h)}\asymp h/\mu$ in this regime.
\end{theorem}

\begin{proof}[Proof of (i)]
Let $e_u:=\ut-u_h$ and $v:=v_\varphi$. Corollary~\ref{cor:GSerr} rearranges to
\[
\frac\mu h\ip{e_u}{v}=\Rc(v)+\Gc(v)-a_h(e_u,v)+\ip{\dn e_u}{v}+\ip{e_u}{\dn v}.
\]
We bound the terms, using that $v$ is smooth up to $O(h^k)$.
\begin{enumerate}
\item $\Rc(v)=G^2+O(\hG^2+h^k)$.
\item $|\Gc(v)|\le C(\hG^{3/2}+\gamma\hG^3+\gamma\hG h^k)$. On the chord, $\ut=\pm d\,\dn u(x^\ast)+O(d^2)$ is tangential while $v\approx-(p-\pbar)n$ is normal. Hence $\ut\cdot v=O(\hG^4+\hG^2h^k)$, the second term coming from $v_\varphi-w$.
\item $|a_h(e_u,v)|\le C_p\norm{\nabla e_u}\le C_pX\le C_p\delta G^2$, by Theorem~\ref{thm:C}.
\item By Theorem~\ref{thm:A} and $h/\mu=\hG/\gamma\le\delta G$,
\[
|\ip{e_u}{\dn v}|\le C_p\norm{e_u}_{\Gh}\le C_p(h/\mu)^{1/2}\tnorm{e_u}\le C_p[(h/\mu)G+X]\le2C_p\delta G^2 .
\]
\item We have
\[
|\ip{\dn e_u}{v}|\le\norm{\dn(\ut-z)}_{\Gh}\norm{v}+C_I(\rho/h)^{1/2}\bigl(\norm{\nabla(z-\ut)}+\norm{\nabla e_u}\bigr)\norm{v}_{\Gh}.
\]
By Theorem~\ref{thm:C} and $(\rho/h)^{1/2}\le(c_0\hG)^{-1/2}$, this is at most $C\hG^{-1/2}XG\le C\delta G^2$.
\end{enumerate}
For $\delta$ small, $|\ip{e_u}{v}|\ge\frac12(h/\mu)G^2$. Dividing by $\norm{v}_{\Gh}\le G+C\hG^2$ gives the lower bound. The upper bound is $\norm{e_u}_{\Gh}\le(h/\mu)^{1/2}\tnorm{e_u}$.
\end{proof}

\begin{proof}[Proof of (ii)]
Let $E\in H^1(\Oh)$ lift $(g-\gI)|_{\partial R}$, supported near $\partial R$, with $E|_{\Gh}=0$ and $\norm{E}_{H^1}\le C\norm{g-\gI}_{H^{1/2}(\partial R)}\le Ch^{k+1/2}$. By Lemma~\ref{lem:korn}, $\norm{e_u}_{\Gh}=\norm{e_u-E}_{\Gh}\le C_{\mathrm{tr}}(\norm{\nabla e_u}+\norm{\nabla E})$.
\end{proof}

\begin{corollary}[Exact leading term]\label{cor:Bp}
Let $G>0$ and $\hG\to0$ with $h/\mu\to0$, $\gamma\hG\to0$ and $\eps/(\hG/\gamma)^{1/2}\to0$. Then
\[
\ut-u_h=(h/\mu)\,v_\varphi+\vartheta_h,\qquad \tnorm{\vartheta_h}=o\bigl((h/\mu)^{1/2}\bigr),\qquad \norm{\vartheta_h}_{\Gh}=o(h/\mu),
\]
and consequently
\[
\frac{\norm{\ut-u_h}_{L^2(\Gh)}}{(h/\mu)G}\to1,\qquad\frac{\tnorm{\ut-u_h}}{(h/\mu)^{1/2}G}\to1 .
\]
\end{corollary}

On $\Gh$, $v_\varphi\approx-(p-\pbar)n$. So to leading order the printed method replaces the no-slip condition by the leak condition $(\mu/h)\,u_h\cdot n\approx p-\pbar$: the penalty, not a traction, balances the pressure.

\begin{proof}
Write $\ut-u_h=(\ut-z)+\delta_R+\delta_G$ in the notation of Theorem~\ref{thm:C}, with $\delta_R=(h/\mu)S_\varphi+r_h$ and $S_\varphi=v_\varphi$.
\begin{enumerate}
\item From step~6 of Theorem~\ref{thm:C},
\[
c_1\tnorm{r_h}^2\le\bigl[C_p(h/\mu+\hG^{3/2})+C\eps\bigr]\tnorm{r_h}.
\]
So $\tnorm{r_h}\le C_p(h/\mu+\hG^{3/2})+C\eps$ and $\norm{r_h}_{\Gh}\le(h/\mu)^{1/2}\tnorm{r_h}$.
\item $\tnorm{\delta_G}+\tnorm{\ut-z}\le C[(1+\gamma^{1/2})\hG^{3/2}+\eps]$, and the $\Gh$ norm carries an extra factor $(h/\mu)^{1/2}$.
\item Relative to $(h/\mu)^{1/2}G$ (energy) and $(h/\mu)G$ (slip), steps 1--2 contribute
\[
O\bigl((h/\mu)^{1/2}/G+(\gamma^{1/2}+\gamma)\hG/G+\eps(\gamma/\hG)^{1/2}/G\bigr)\to0 .
\]
\item By Lemma~\ref{lem:chords} (Jacobian and displacement of $\Gh$ from $\Gamma$), $\norm{v_\varphi}_{\Gh}=G(1+O(\hG^2))+O(h^k)$ and $\tnorm{v_\varphi}^2=(\mu/h)\norm{v_\varphi}^2_{\Gh}+O(1)$.\qedhere
\end{enumerate}
\end{proof}

This gives the sharp answer to Scott's question for the printed method. The estimate $\hG^{3/2}+h^k$ fails exactly when $p$ is not constant on $\Gamma$. The error is then the explicit field $(h/\mu)v_\varphi$, with the universal constant $1$ in both the slip norm and the energy norm.

\subsection{The \texorpdfstring{$H^1$}{H1} leak constant}\label{hc:sec}

Corollary~\ref{cor:Bp} identifies the slip and energy constants of $\GS$ (both $\to1$), but not the $H^1$ constant
$\norm{\nabla(\ut-u_h)}\,\mu/(hG)$, because the remainder $r_h=\delta_R-(h/\mu)S_\varphi$ of Theorem~\ref{thm:C} has the same
$H^1$ order as $(h/\mu)S_\varphi$. We identify the limit of $(\mu/h)\delta_R$: it is the Stokes flow in $\Omega$ driven by the
leak as \emph{Dirichlet} data. The $H^1$ constant is $K=\norm{\nabla U}/G$, a global quantity that depends on the outer
box (it tends to $3/\sqrt7$ as $L\to\infty$ for the pressure of test P). For test P and $L=2.5$ a converged series computation gives
$K=2.7565157$, inside the rigorous bracket $[2.0746,\,3.0526]$ of Proposition~\ref{hc:bracket}.

\subsubsection{The limit problem}

\begin{definition}[Leak flow]\label{hc:def}
Let $(U,P)$ solve
\[
-\Delta U+\nabla P=0,\quad \div U=0\ \text{in }\Omega,\qquad U=(p-\pbar)\,e_r=-(p-\pbar)\,n\ \text{on }\Gamma,\qquad U=0\ \text{on }\partial R ,
\]
and put $K:=\norm{\nabla U}_{L^2(\Omega)}/G$.
\end{definition}

The data are compatible, since $\int_\Gamma U\cdot n=-\int_\Gamma(p-\pbar)=0$. Hence $(U,P)$ exists and is unique ($P$ up to
a constant), $U=\curl\Psi$ with a single-valued stream function, and $U|_\Gamma=w|_\Gamma$ with $w=\curl\varphi$ the test field
above. $U$ is smooth up to $\Gamma$ (smooth boundary and data) and $U\in H^2(\Omega)$ (the corners of $R$
are convex). Since $-\Delta U=-\nabla P$, for every divergence-free $V$ with the same boundary values
$\int\nabla U:\nabla(V-U)=-\int\nabla P\cdot(V-U)=0$, so
\begin{equation}\label{hc:min}
\norm{\nabla U}^2=\min\Bigl\{\norm{\nabla V}^2_{L^2(\Omega)}:\ V\in H^1(\Omega)^2,\ \div V=0,\ V|_\Gamma=(p-\pbar)e_r,\ V|_{\partial R}=0\Bigr\}.
\end{equation}
In particular $K$ depends on $p|_\Gamma$ \emph{and} on the outer boundary: it is not a local quantity.

\subsubsection{Convergence}

Extend $\Psi$ smoothly across $\Gamma$ to a fixed neighbourhood and put $\tilde U:=\curl\tilde\Psi$ on $\Oh$, so $\div\tilde U=0$;
extend $P$ smoothly to $\tilde P$. Then $F_U:=-\Delta\tilde U+\nabla\tilde P$ vanishes on $\Omega$, is bounded, and is supported in
the strip $S_\varphi$.

\begin{lemma}\label{hc:approx}
There is $z_U\in\Zh$ with $\tnorm{\tilde U-z_U}\le E_U:=C\bigl(h\norm{U}_{H^2(\Omega)}+(1+(\gamma\hG)^{1/2})h^k+(1+\gamma^{1/2})\hG^{\sigma_k-1/2}\bigr)$. For bounded $\gamma$, $E_U\le Ch$, since $\sigma_k-\frac12\ge k+\frac12$.
\end{lemma}

\begin{proof}
Repeat the proof of Lemma~\ref{lem:approx} with $\ut$ replaced by $\tilde U$ and $g=0$. Away from $\Gamma$ only $\tilde U\in
H^2$ is available, so the interpolation step gives $\norm{\nabla(\tilde U-I_h\tilde U)}\le Ch\norm{U}_{H^2(\Omega)}$. Near $\Gamma$, $\tilde U$
is smooth and the boundary terms are as before. The net flux $m=\int_{\Gh}(I_h\tilde U-\tilde U)\cdot n$ is $O(\hG^{\sigma_k})$
because $\int_{\Gh}\tilde U\cdot n=\int_{\Oh}\div\tilde U=0$. As in step~3 of that proof, the flux bubble then costs
$C|m|(\hG^{-1/2}+(\mu/h)^{1/2})\le C(1+\gamma^{1/2})\hG^{\sigma_k-1/2}$; the Bogovskii correction (Lemma~\ref{lem:infsup})
is unchanged.
\end{proof}

\begin{theorem}[The leak flow is the $H^1$ limit]\label{hc:thm}
Let $\gamma\ge\gamma_0$ and let $\delta_R\in\Zh$ solve $N_h(\delta_R,v)=\Rc(v)$ for all $v\in\Zh$ (Theorem~\ref{thm:C}, step 1).
Then
\[
\Bigl\|\nabla\Bigl(\frac\mu h\delta_R-\tilde U\Bigr)\Bigr\|_{\Oh}
\le\tnorm{\frac\mu h\delta_R-z_U}+E_U
\le C_U\Bigl[\min\bigl(\gamma\hG^{1/2},(\gamma\hG)^{1/2}\bigr)+\hG^{1/2}+(h/\mu)^{1/2}+E_U\Bigr],
\]
where $C_U$ depends on (M0)--(M2), $k$, $L$ and on norms of $p$ and $U$ near $\Gamma$.
\end{theorem}

\begin{proof}
Put $\epsilon:=\delta_R-(h/\mu)z_U\in\Zh$. For $v\in\Zh$,
\begin{equation}\label{hc:erreq}
N_h(\epsilon,v)=\Rc(v)-\frac h\mu N_h(\tilde U,v)+\frac h\mu N_h(\tilde U-z_U,v).
\end{equation}
\begin{enumerate}
\item \emph{Green's formula for $\tilde U$.} As in Lemma~\ref{lem:erroreq}, $\div D(\tilde U)=\Delta\tilde U$, $\div v=0$, $v|_{\partial R}=0$ and
$\int_{\Gh}v\cdot n=0$ give, for any constant $c$,
\[
N_h(\tilde U,v)=(F_U,v)+\ip{\sigma_U}{v}-\ip{\tilde U}{\dn v}+\frac\mu h\ip{\tilde U}{v},\qquad
\sigma_U:=(\nabla\tilde U)^{\mathsf T}n-(\tilde P-c)\,n .
\]
\item \emph{The penalty terms cancel to leading order.} With $\Rc(v)=-\ip{(\pt-\pbar)n}{v}$, \eqref{hc:erreq} becomes
\[
N_h(\epsilon,v)=-\ip{\zeta}{v}-\frac h\mu\Bigl[(F_U,v)+\ip{\sigma_U}{v}-\ip{\tilde U}{\dn v}\Bigr]+\frac h\mu N_h(\tilde U-z_U,v),
\qquad \zeta:=(\pt-\pbar)n+\tilde U .
\]
\item \emph{The mismatch $\zeta$.} On $e\in\EG$, by Lemma~\ref{lem:chords} and $\tilde U(x^\ast)=-(p-\pbar)(x^\ast)n(x^\ast)$,
\[
\zeta=(p-\pbar)(x^\ast)\bigl(n_e-n(x^\ast)\bigr)+O(\hG^2)=(p-\pbar)(m_e^\ast)\,s\,\tau_e+O(\hG^2),
\]
exactly the functional $m$ of Theorem~\ref{thm:C}, step~4 (without the $h^k$ term). Hence
$|\ip{\zeta}{v}|\le C\hG^{3/2}\norm{\nabla v}$. Alternatively $\norm{\zeta}_{\Gh}\le C\hG$ gives
$|\ip{\zeta}{v}|\le C\hG(h/\mu)^{1/2}\tnorm v$.
\item \emph{The $O(h/\mu)$ terms.} $|(F_U,v)|\le C\hG^2\norm{\nabla v}$ (Lemma~\ref{lem:ext}); $|\ip{\sigma_U}{v}|\le
C\norm{v}_{\Gh}\le C(h/\mu)^{1/2}\tnorm v$. On $\Gh$, $\tilde U$ is edgewise smooth, continuous at the vertices, and
$\tilde U\cdot\tau_e=-(p-\pbar)(x^\ast)\,n(x^\ast)\cdot\tau_e+O(\hG^2)=O(\hG)$. So Lemma~\ref{lem:cancel} applies with
$S=\tilde U$: $|\ip{\tilde U}{\dn v}|\le C(\norm v_{\Gh}+\hG^{1/2}\norm{\nabla v})\le C((h/\mu)^{1/2}+\hG^{1/2})\tnorm v$.
Finally $|N_h(\tilde U-z_U,v)|\le ME_U\tnorm v$ (Lemmas~\ref{lem:coercive} and~\ref{hc:approx}).
\item \emph{Close.} Test with $v=\epsilon$ and use Lemma~\ref{lem:coercive}:
\[
c_1\tnorm{\epsilon}\le C\Bigl[\min\bigl(\hG^{3/2},\hG(h/\mu)^{1/2}\bigr)+\frac h\mu\bigl(\hG^2+(h/\mu)^{1/2}+\hG^{1/2}+E_U\bigr)\Bigr].
\]
Divide by $h/\mu$, and use $\hG^{3/2}\mu/h=\gamma\hG^{1/2}$ and $\hG(\mu/h)^{1/2}=(\gamma\hG)^{1/2}$. The first inequality of the
theorem is the triangle inequality with Lemma~\ref{hc:approx}.\qedhere
\end{enumerate}
\end{proof}

\begin{corollary}[The $H^1$ constant]\label{hc:cor}
\begin{enumerate}[label=(\roman*)]
\item For $\GS$ with general data, in the notation of Theorem~\ref{thm:C},
\[
\Bigl\|\nabla\Bigl(\frac\mu h(\ut-u_h)-\tilde U\Bigr)\Bigr\|_{\Oh}\le\text{(bound of Theorem~\ref{hc:thm})}
+C\frac\mu h\Bigl[(1+\gamma^{1/2})\hG^{3/2}+\eps\Bigr].
\]
\item Consequently, if $\hG\to0$ with $h/\mu\to0$, $\gamma\hG\to0$, $\gamma(1+\gamma^{1/2})\hG^{1/2}\to0$ and
$(\mu/h)\eps\to0$ (all of which hold for fixed $\gamma$ and bounded $\rho$), then
\[
\frac{\norm{\nabla(\ut-u_h)}_{\Oh}}{(h/\mu)\,G}\longrightarrow K=\frac{\norm{\nabla U}_{L^2(\Omega)}}{G}.
\]
\item On test P ($u=0$, $f=\nabla p$) one has $\ut=0$, $z=0$ and $\Gc\equiv0$, so $\ut-u_h=\delta_R$ exactly and
Theorem~\ref{hc:thm} applies with no further condition.
\item The remainder of Theorem~\ref{thm:C} converges: $(\mu/h)\,r_h\to U-w$ in $H^1$. The limit is the Stokes correction
of the cut-off field $w$: it vanishes on $\partial\Omega$ and solves $-\Delta(U-w)+\nabla P=\Delta w$. It depends on the
cut-off $\chi_1$, whereas $U$ does not.
\end{enumerate}
\end{corollary}

\begin{proof}
(i) is $\ut-u_h=(\ut-z)+\delta_R+\delta_G$ with steps 2 of Theorem~\ref{thm:C} and Lemma~\ref{lem:approx}. (ii) Multiply out;
$\norm{\nabla\tilde U}^2_{\Oh}-\norm{\nabla U}^2_{\Omega}=\norm{\nabla\tilde U}^2_{S_h}=O(\hG^2)$. (iii) $F=\ft-\nabla\pt=0$ and all terms of
$\Gc$ contain $\ut$. (iv) $r_h=\delta_R-(h/\mu)v_\varphi$ and $\norm{\nabla(v_\varphi-w)}\le Ch^k$.
\end{proof}

\begin{remark}[The Dirichlet-dominated regime]\label{hc:regime}
The hypothesis $\gamma\hG=\mu\hG^2/h\to0$ is the one of Corollary~\ref{cor:Bp}. It fails for the large-$\mu$ runs of
Table~\ref{tab:P} ($\gamma\hG\approx240$ at $N=96$, $\mu=10^4$). There the discrete problem is close to its $\mu\to\infty$ limit,
a discrete Dirichlet problem whose data $-(p-\pbar)n_e$ jumps in direction at every vertex of $\Gh$. The computations below
show that $(\mu/h)\delta_R\to U$ in that regime too, with relative $H^1$ distance $d_h$ decaying roughly like $\hG^{0.6}$--$\hG^{0.8}$, the decay slowing as $N$ grows, but
Theorem~\ref{hc:thm} does not cover it. Proving convergence uniformly in $\mu\ge\mu_0$ is open.
\end{remark}

\subsubsection{Rigorous bracket and exact asymptotics}

For $\varrho>1$ let $\mathcal A_\varrho=\{1<r<\varrho\}$ and let $\mathcal E(\varrho)$ be the minimum \eqref{hc:min} over $\mathcal A_\varrho$ (zero data on
$r=\varrho$). Extending by zero shows that the minimum decreases as the domain grows.

\begin{proposition}\label{hc:bracket}
If $\mathcal A_{\varrho_1}\subset\Omega\subset\mathcal A_{\varrho_2}$, then $\mathcal E(\varrho_2)\le\norm{\nabla U}^2\le\mathcal E(\varrho_1)$. In the annulus the
Fourier modes decouple: if $\Psi|_\Gamma=\sum_n(a_n\cos n\theta+b_n\sin n\theta)$, then $\mathcal E(\varrho)=\sum_n(a_n^2+b_n^2)E_n(\varrho)$, where
$E_n(\varrho)=\pi\int_1^\varrho\bigl[f''^2+2n^2((f/r)')^2+(f'/r-n^2f/r^2)^2\bigr]r\,dr$, and $f$ is the combination of
$r^n,r^{-n},r^{n+2},r^{2-n}$ ($r,r^{-1},r^3,r\log r$ for $n=1$) with $f(1)=1$, $f'(1)=0$, $f(\varrho)=f'(\varrho)=0$.
\end{proposition}

For test P, $p-\pbar=-\frac34\sin\theta+\frac14\sin3\theta+\frac12\cos\theta$ on $\Gamma$, so
$\Psi|_\Gamma=\frac34\cos\theta+\frac12\sin\theta-\frac1{12}\cos3\theta$ and $G^2=7\pi/8$. With $\varrho_1=L=2.5$ and $\varrho_2=L\sqrt2$
(40-digit quadrature):
\[
\boxed{\,2.0746004\le K\le3.0525769\,}
\]
The same formulas give the dependence on the box. $E_3(\varrho)\to45\pi$ (the exterior clamped solution is
$f=\frac32r^{-1}-\frac12r^{-3}$). $E_1(\varrho)\to\pi$: the $n=1$ data split into a potential dipole (exterior energy exactly $\pi$) and
a rigid translation, whose energy vanishes as $\varrho\to\infty$ (the Stokes paradox). The annulus values $E_1(\varrho)-\pi=2.47$, $0.87$, $0.38$ at $\varrho=10,10^2,10^4$ are consistent with the classical $1/\log\varrho$ rate, which we observe numerically here rather than prove. Hence
\[
K(L)\longrightarrow\sqrt{\tfrac{(\frac9{16}+\frac14)\pi+\frac{45\pi}{144}}{7\pi/8}}=\frac3{\sqrt7}\approx1.134\qquad(L\to\infty),
\]
but only logarithmically: the annulus values are $K=1.420$, $1.242$, $1.183$ at $\varrho=10,10^2,10^4$. So the constant
$2.7$--$2.8$ of Table~\ref{tab:P} belongs to the box $L=2.5$, not to the method.

\subsubsection{The value for test P}

\emph{Series solution.} Write $\Psi=\sum_{n\text{ odd}\le n_{\max}}f_n(r)\cos n\theta$ for data $\cos\theta$ and $\cos3\theta$ separately (the class of
$\cos n\theta$, $n$ odd, is invariant under both reflections of the square; the $\sin\theta$ datum is the rotation of the
$\cos\theta$ one by $\pi/2$, and cross terms between the two reflection classes vanish). Imposing $\Psi=0$ on $\partial R$, rather than an unknown constant, is justified by symmetry: for data $\cos n\theta$ with $n$ odd, the reflection $x\mapsto-x$ maps $\Omega$ to itself and changes the sign of the data, so by uniqueness $\Psi$ is odd under it, and its constant value on $\partial R$ is zero; the $\sin\theta$ datum follows by rotation. The conditions on $r=1$ are imposed
exactly, mode by mode, and $\Psi=\partial_\nu\Psi=0$ on $\partial R$ by least squares at $16\,000$ Chebyshev-clustered points.
The energy is integrated by Gauss quadrature on eight polar sectors. The value is stable to 8 digits under refinement of both the truncation $n_{\max}$ and the
quadrature:

\begin{center}
\begin{tabular}{rccc}
\toprule
$n_{\max}$ & max boundary residual & $\norm{\nabla U}$ & $K$\\
\midrule
15 & $4.4\cdot10^{-3}$ & 4.57012143 & 2.75644148\\
31 & $6.1\cdot10^{-5}$ & 4.57024239 & 2.75651444\\
47 & $5.1\cdot10^{-6}$ & 4.57024449 & 2.75651571\\
63 & $1.1\cdot10^{-6}$ & 4.57024450 & 2.75651571\\
79 & $1.8\cdot10^{-7}$ & 4.57024450 & 2.75651571\\
\bottomrule
\end{tabular}
\end{center}

\[
\boxed{\,K_{\mathrm P}=2.7565157\,}
\]
This lies inside the rigorous bracket.
\cert{code/h1\_limit.py annulus; code/h1\_limit.py mps}{logs/h1\_limit\_series.log}

\emph{Finite element check.} Our solver (SuperLU with refinement, as in Section~\ref{sec:numerics}) was run on test P, and
the discrete field was compared \emph{pointwise} with the series $U$ at the quadrature points. Both the normalised $H^1$ error
$K_h$ and the relative distance $d_h:=\norm{\nabla((\mu/h)(\ut-u_h)-\tilde U)}/\norm{\nabla U}$ are shown below. The $\mu=100$ and
$N=96$ values of $K_h$ reproduce Table~\ref{tab:P} to all printed digits.

\begin{center}\small
\begin{tabular}{r|cc|cc|cc|cc}
\toprule
 & \multicolumn{2}{c|}{$\mu=10^2$} & \multicolumn{2}{c|}{$\mu=10^3$} & \multicolumn{2}{c|}{$\mu=10^4$} & \multicolumn{2}{c}{$\mu=10^8$}\\
$N$ & $K_h$ & $d_h$ & $K_h$ & $d_h$ & $K_h$ & $d_h$ & $K_h$ & $d_h$\\
\midrule
16 & 2.4642 & 0.193 & 2.7844 & 0.318 & 3.0562 & 0.539 & 3.3254 & 0.717\\
24 & 2.5629 & 0.147 & 2.7913 & 0.259 & 2.9309 & 0.398 & 3.0452 & 0.497\\
32 & 2.6117 & 0.124 & 2.7880 & 0.222 & 2.8751 & 0.323 & 2.9339 & 0.385\\
48 & 2.6601 & 0.098 & 2.7808 & 0.179 & 2.8283 & 0.246 & 2.8514 & 0.277\\
64 & 2.6843 & 0.084 & 2.7759 & 0.154 & 2.8082 & 0.206 & 2.8208 & 0.226\\
96 & 2.7084 & 0.067 & 2.7701 & 0.124 & 2.7899 & 0.163 & 2.7959 & 0.175\\
\bottomrule
\end{tabular}
\end{center}
\cert{code/h1\_limit.py fem}{logs/h1\_limit\_fem.log}

The measured values $2.71$--$2.79$ bracket the limit $2.7565$ from both sides. The approach runs in two directions.
\begin{itemize}
\item At $\mu=100$, $K_h$ increases towards $K$ and $d_h$ decays like $h^{0.6}$ ($0.193\to0.067$ as $h$ drops by a factor $5.3$).
This is consistent with Theorem~\ref{hc:thm}: for fixed $\mu$, the terms $(h/\mu)^{1/2}$ and $\hG^{1/2}$ are both
$\propto h^{1/2}$, and $(\gamma\hG)^{1/2}\propto h^{1/2}$ too.
\item At large $\mu$ (Remark~\ref{hc:regime}), $K_h$ decreases towards $K$ from above. The excess is the energy of an
element-scale oscillation that is nearly orthogonal to $U$: $K_h^2\approx K^2(1+d_h^2)$, for example $2.7565\sqrt{1+0.175^2}=2.798$
against $2.796$ at $N=96$.
\end{itemize}
At fixed $N=96$ the $\mu$-dependence of Table~\ref{tab:P} ($2.708\to2.770\to2.790$) is the crossover between the two regimes,
not a sign of a different limit.
